\section{Theta characteristics and a complete double cover}
\label{sec:theta}

The fixed opposite-distance difference $1/5$ gives a particularly explicit
connection between the distance equations, the geometry of a plane quartic,
and descent on its Jacobian.  In this section we determine all the odd theta
characteristics of that quartic and identify the double cover which restores
the four distances.  A rational point at infinity then gives a local
solubility theorem for every finite multiplication cover, including the
pullbacks which retain this double cover.

\subsection{The quartic quotient and its metric interpretation}

Let $C\subset\PP^4_{[P:Q:R:z:T]}$ be the curve
\begin{equation}
 \begin{split}
 P^2&=z^2+50T^2,\\
 Q^2&=z^2+14zT+51T^2,\\
 R^2&=z^2-14zT+51T^2.
 \end{split}
 \label{eq:theta-C}
\end{equation}
We use $C$ as a complete source: all three equations are imposed at the same
value of $z/T$.  On $T=1$, put
\begin{equation}
 \begin{aligned}
 x&=\frac{50-P-7z}{100},&
 y&=\frac{50-P+7z}{100},\\
 (r_0,r_1,r_2,r_3)&=
 \left(\frac{P-1}{10},\frac Q{10},\frac R{10},\frac{P+1}{10}\right).
 \end{aligned}
 \label{eq:theta-return}
\end{equation}
The strict positive locus $\mathcal U\subset C(\R)$ is defined by
\begin{equation}
 T=1,\qquad P>1,\quad Q>0,\quad R>0,\qquad P+7|z|<50.
 \label{eq:theta-U}
\end{equation}
It corresponds exactly to the strict distance configurations with
$r_3-r_0=1/5$.  For completeness, the first distance identity follows from
\[
 5000\bigl(r_0^2-x^2-y^2\bigr)
   =49(P^2-z^2-50)=0.
\]
The other three follow from
\[
 r_1^2-r_0^2=1-2x,\qquad
 r_2^2-r_0^2=1-2y,\qquad
 r_3^2-r_0^2=2-2x-2y.
\]
Conversely, the inverse formulas are
\[
 P=5(r_3+r_0),\quad Q=10r_1,\quad R=10r_2,
 \quad z=\frac{50(y-x)}7.
\]
The inequalities in \eqref{eq:theta-U} are precisely the positivity of the
chosen roots and $0<x,y<1$.

\begin{proposition}
\label{prop:theta-quotient}
The curve $C$ is smooth of genus five and has eight rational points at
$T=0$.  The fixed-point-free involution
\[
 \sigma[P:Q:R:z:T]=[-P:-Q:-R:z:T]
\]
has quotient the smooth plane quartic $D\subset\PP^2_{[X:Y:Z]}$ defined by
\begin{equation}
 F=9801(X^4+Y^4)+19598X^2Y^2
       -39592(X^2+Y^2)Z^2+39984Z^4=0.
 \label{eq:theta-quartic}
\end{equation}
The quotient map is
\[
 \pi:C\longrightarrow D,\qquad [P:Q:R:z:T]\longmapsto[Q:R:P].
\]
In particular, $\pi$ is a connected \emph{\'etale} double cover.  The points
\[
 O_C=[1:1:1:1:0],\qquad O_D=[1:1:1]
\]
are rational and satisfy $\pi(O_C)=O_D$.
\end{proposition}

\begin{proof}
On the chart $T=1$, the radicands in \eqref{eq:theta-C} have
discriminants $-200$, $-8$ and $-8$, and pairwise resultants
$9801$, $9801$ and $39984$. Their six roots are therefore distinct.  A simple zero
of each polynomial shows that its square class is independent of those of
the other two over $\overline\Q(z)$.  The map to the $[z:T]$-line has
degree eight, with six branch values and inertia order two at each.
Riemann--Hurwitz gives
\[
 2g(C)-2=8(-2)+6\cdot4=8.
\]
The simple branch points are smooth.  At $T=0$ we may set $z=1$, and the
points are exactly
\[
 [\epsilon_1:\epsilon_2:\epsilon_3:1:0],\qquad
 \epsilon_i\in\{1,-1\}.
\]
The three derivatives with respect to $P,Q,R$ are nonzero there, so these
points are smooth as well.

There is no point of $C$ with $P=Q=R=0$: subtracting the last two equations
gives $zT=0$, and the equations then force $z=T=0$.  Hence $\pi$ is a
morphism.  A projective fixed point of $\sigma$ would have either
$P=Q=R=0$ or $z=T=0$, both impossible.

On the chart $P\ne0$, set
\[
 \xi=Q/P,\qquad \eta=R/P,\qquad v=T/P,\qquad w=z/P,
 \qquad S=\xi^2+\eta^2.
\]
The three quadrics give
\begin{equation}
 v^2=\frac{S-2}{2},\qquad
 wv=\frac{\xi^2-\eta^2}{28},\qquad
 w^2=51-25S.
 \label{eq:theta-vw}
\end{equation}
Eliminating $v,w$ gives
\begin{equation}
 (\xi^2-\eta^2)^2+9800(S-2)^2-392(S-2)=0,
 \label{eq:theta-elimination}
\end{equation}
which is \eqref{eq:theta-quartic} on $Z=1$.

Here is also a direct smoothness proof for the plane equation.  Write
\[
 a=9801,\quad b=19598,\quad c=-39592,\quad d=39984,
 \qquad
 M=\begin{pmatrix}2a&b&c\\b&2a&c\\c&c&2d\end{pmatrix}.
\]
At a singular point, the rows indexed by the nonzero coordinates
annihilate the corresponding vector of coordinate squares.  The nonempty
principal minors of $M$ have values
\[
 19602,\quad79968,\quad156800,\quad6272,\quad-1229312,
\]
with repetitions.  Every such submatrix is invertible.  This proves
smoothness; a smooth plane quartic is geometrically integral and has
genus three.  The map to it has degree two, since generically $v$ determines
$w$ in \eqref{eq:theta-vw}, and the two choices are interchanged by $\sigma$.
The fixed-point calculation proves that the map is \emph{\'etale}.
\end{proof}

\begin{proposition}
\label{prop:theta-strict-lift}
A rational point $[\xi:\eta:1]\in D(\Q)$ is the image of a point of
$\mathcal U\cap C(\Q)$ if and only if
\begin{equation}
 \xi>0,\qquad\eta>0,\qquad
 v=\sqrt{\frac{\xi^2+\eta^2-2}{2}}\in\Q_{>0},\qquad
 \left|\frac{\xi^2-\eta^2}{28v}\right|<\frac7{11}.
 \label{eq:theta-strict-lift}
\end{equation}
For this choice of $v$, the unique point on the positive branch is
\begin{equation}
 P=\frac1v,\qquad Q=\frac\xi v,\qquad R=\frac\eta v,
 \qquad z=\frac wv,
 \qquad w=\frac{\xi^2-\eta^2}{28v}.
 \label{eq:theta-inverse}
\end{equation}
\end{proposition}

\begin{proof}
Equation \eqref{eq:theta-elimination} gives $w^2+50v^2=1$, and substitution
in \eqref{eq:theta-inverse} gives all three equations of $C$.
Moreover $P\ge\sqrt{50}>1$.  The remaining strict inequality is
\[
 P+7|z|<50\quad\Longleftrightarrow\quad1+7|w|<50v.
\]
Both sides on the right are positive.  Squaring and using
$50v^2=1-w^2$ reduces it to
\[
 99|w|^2+14|w|-49<0,
\]
whose positive root is $7/11$.  The inverse implication follows from the
same calculation.
\end{proof}

Every real point of $D$ lies on $Z\ne0$.  Indeed,
\[
 F(X,Y,0)=9801(X^2+Y^2)^2-4X^2Y^2>0
\]
for a nonzero real pair $(X,Y)$.  The four points where $v^2=0$ are
\begin{equation}
 E_{\epsilon,\eta}=[\epsilon:\eta:1],\qquad
 \epsilon,\eta\in\{1,-1\};
 \label{eq:theta-four-boundary}
\end{equation}
each has two rational preimages at $T=0$ on $C$.

\subsection{The complete bitangent configuration}

Put $r=\sqrt2$, $s=\sqrt{51}$, and $K=\Q(r,s)$.  Independent signs in the
following statement range over $\{1,-1\}$.

\begin{theorem}
\label{thm:theta-lines}
The twenty-eight bitangents of $D$ are precisely the following lines:
\begin{enumerate}
\item $Z=uX$ and $Z=uY$, where $u=\epsilon r/2$ or $u=5\epsilon r$;
\item $Z=(\epsilon hX+\eta hY)/14$, where $h=9$ or $h=11$;
\item $Z=s(\epsilon hX+\eta jY)/102$, where
      $(h,j)=(9,11)$ or $(11,9)$;
\item $X=tY$, where $t=\epsilon r(s+7\eta)/2$.
\end{enumerate}
Each has two distinct contact points.  The associated odd-theta
\emph{\'etale} algebra is
\begin{equation}
 L_{\mathrm{odd}}
 =\Q^8\times\Q(\sqrt2)^4\times\Q(\sqrt{51})^4
          \times\Q(\sqrt2,\sqrt{51}).
 \label{eq:theta-algebra}
\end{equation}
The full Galois action is induced by the independent sign changes of $r$
and $s$ in the displayed lines.
\end{theorem}

\begin{proof}
We first give the contact identities.  These also specify the contact
divisors, since each is the zero divisor of the displayed quadratic on its
line.  For type~1, on $Z=uX$ one has
\begin{equation}
 F|_{Z=uX}=9801(Y^2+\alpha X^2)^2,
 \quad
 \alpha=\begin{cases}
 -1/99,&u^2=1/2,\\
 -9899/99,&u^2=50.
 \end{cases}
 \label{eq:theta-contact-I}
\end{equation}
Exchange $X,Y$ for the other four lines.  For type~2,
\begin{equation}
 \begin{array}{c|c|c}
 h&k_h&q_h\\\hline
 9&-818/81&13122/49\\
 11&1222/121&29282/49
 \end{array}
 \qquad
 F|_{Z=h(\epsilon X+\eta Y)/14}
   =q_h(X^2+\epsilon\eta k_hXY+Y^2)^2.
 \label{eq:theta-contact-II}
\end{equation}
For type~3 the identity is
\begin{equation}
 F|_{Z=s(\epsilon hX+\eta jY)/102}
   =\frac{6534}{17}(X^2+\epsilon\eta k_{h,j}XY-Y^2)^2,
 \quad
 k_{9,11}=-\frac{980}{99},\quad k_{11,9}=\frac{980}{99}.
 \label{eq:theta-contact-III}
\end{equation}
Finally, $t^2=50+7\eta s$ for type~4, and
\begin{equation}
 F|_{X=tY}
  =39984\left(Z^2-\frac{39592(t^2+1)}{79968}Y^2\right)^2.
 \label{eq:theta-contact-IV}
\end{equation}
All these homogeneous quadratics have nonzero discriminant: in
\eqref{eq:theta-contact-I} the coefficient $\alpha$ is nonzero; in
\eqref{eq:theta-contact-II} one has $k_h\ne\pm2$; in
\eqref{eq:theta-contact-III} one has $k_{h,j}^2+4>0$; and in
\eqref{eq:theta-contact-IV} one has $t^2+1\ne0$.

We prove completeness by comparing coefficients of binary quartics.
A line with one vanishing coefficient is of the form $X=a_0Z$,
$Y=a_0Z$, or $X=a_0Y$.  For the first two forms the discriminant condition
for the restricted even quartic to be a square is
\[
 -3136(a_0^2-2)(50a_0^2-1)=0.
\]
For the third form it is
\[
 -6272(a_0^4-100a_0^2+1)=0.
\]
The resulting lines are exactly types~1 and~4.  The three coordinate axes
give nonzero discriminants and hence contribute no line.

Now write a line with all coefficients nonzero as $Z=\alpha X+\beta Y$,
and set $U=\alpha^2$, $V=\beta^2$.  With $a,b,c,d$ as in the proof of
Proposition~\ref{prop:theta-quotient}, its restricted quartic has
coefficients
\begin{align*}
 A_4&=a+cU+dU^2,&
 A_3&=2\alpha\beta(c+2dU),\\
 A_2&=b+c(U+V)+6dUV,&
 A_1&=2\alpha\beta(c+2dV),\\
 A_0&=a+cV+dV^2.
\end{align*}
For a square quartic, $A_4=0$ would imply $A_3=0$.  The latter forces
$U=101/204$, for which $A_4=2/51$.  Thus $A_4,A_0$ are nonzero.
If one of $A_3,A_1$ vanishes, both vanish, giving $U=V=101/204$; at that
pair
\[
 A_2^2-4A_4A_0=\frac{78368326400}{51}\ne0.
\]
In the remaining case, comparison with the square of a quadratic gives
\begin{align}
 &(c+2dU)^2A_0-(c+2dV)^2A_4\notag\\
 &\hspace{22mm}=2458624(U-V)(102U+102V-101)=0,
 \label{eq:theta-square-factor-I}\\
 E&:=A_0(c+2dV)A_2-UV(c+2dV)^3-2A_0^2(c+2dU)=0.
 \label{eq:theta-square-factor-II}
\end{align}
More explicitly, when these four coefficients are nonzero, the restricted
quartic is a square precisely when
\[
 A_4=A_0(A_3/A_1)^2,\qquad
 A_2=\frac{A_1^2}{4A_0}+\frac{2A_0A_3}{A_1}.
\]
The square is then
\[
 A_0\left(\frac{A_3}{A_1}X^2+
                   \frac{A_1}{2A_0}XY+Y^2\right)^2.
\]
Substitution gives exactly
\eqref{eq:theta-square-factor-I}--\eqref{eq:theta-square-factor-II}.
On the two lines singled out by
\eqref{eq:theta-square-factor-I}, the second expression factors as
\begin{align*}
 E|_{V=U}
  &=-1568(196U-121)(196U-81)(204U-101),\\
 E|_{V=101/102-U}
  &=\frac{76832}{17}(68U-27)(204U-121)(204U-101).
\end{align*}
After discarding $U=101/204$, already treated, these yield
\[
 (U,V)=\left(\frac{81}{196},\frac{81}{196}\right),\quad
       \left(\frac{121}{196},\frac{121}{196}\right),\quad
       \left(\frac{27}{68},\frac{121}{204}\right),\quad
       \left(\frac{121}{204},\frac{27}{68}\right).
\]
These are precisely types~2 and~3, and their sufficiency was established
by \eqref{eq:theta-contact-II}--\eqref{eq:theta-contact-III}.

Since $D$ is a smooth plane quartic, its canonical bundle is
$\OO_D(1)$.  A bitangent $\ell$ therefore cuts out $2B_\ell$, where
$B_\ell$ is an effective theta characteristic of degree two.
It has $h^0(B_\ell)=1$: a degree-two pencil would make $D$
hyperelliptic, whereas its canonical map is a plane embedding.
Conversely, an odd theta characteristic has a unique effective divisor
of degree two and gives a bitangent by taking twice that divisor.
The preceding classification is thus the full odd-theta set.

The eight type~2 lines are rational.  Type~1 consists of four quadratic
orbits over $\Q(r)$, and type~3 consists of four quadratic orbits over
$\Q(s)$.  The four type~4 lines form one orbit over $K$: the value of
$t^2$ recovers $s$ up to sign, and $t$ then recovers $r$ up to sign.
This proves \eqref{eq:theta-algebra} and the Galois assertion.
\end{proof}

\subsection{Theta fractions and the literal lifting class}

Choose the rational bitangent
\[
 \ell_0=14Z-9X-9Y.
\]
None of the twenty-eight bitangents contains $O_D$: this follows directly
from their equations (for type~3 the coefficient of $s$ is rational and
nonzero, and for type~4 $t^2=50\pm7s\ne1$).  Thus, for every line in
Theorem~\ref{thm:theta-lines}, the function
\begin{equation}
 f_\ell=\frac{\ell}{\ell_0}\frac{\ell_0(O_D)}{\ell(O_D)}
 \label{eq:theta-fractions}
\end{equation}
is defined at $O_D$ and satisfies
\[
 f_\ell(O_D)=1,\qquad
 \operatorname{div}(f_\ell)=2(B_\ell-B_{\ell_0}).
\]
The line equations and their Galois action therefore give explicitly
labelled theta fractions in each factor of \eqref{eq:theta-algebra}.
Alternatively, relative to the canonical divisor $H=(Z=0)|_D$, the
functions $\ell/Z$ have divisor $2B_\ell-H$.

\begin{proposition}
\label{prop:theta-cover-class}
Let
\[
 C_0=X^2+Y^2-2Z^2,\qquad
 N=97X^2+99Y^2-196Z^2,\qquad
 h=\frac{C_0}{2Z^2},\qquad
 \ell_\pm=Z\pm\frac X{\sqrt2}.
\]
The cover $\pi:C\to D$ is the normalization of $v^2=h$.  If
$E_0=\sum_{\epsilon,\eta}E_{\epsilon,\eta}$ is the sum of the four points
in \eqref{eq:theta-four-boundary}, its geometric covering class is
\[
 \delta=[E_0-H]\in\Jac(D)(\Q)[2]\setminus\{0\}.
\]
Over $\Q(\sqrt2)$ it is identified with
$[B_{\ell_+}-B_{\ell_-}]$ by the exact function-field identity
\begin{equation}
 h=\frac{\ell_+}{\ell_-}
       \left(\frac{N}{28\sqrt2\,Z\ell_+}\right)^2.
 \label{eq:theta-cover-identity}
\end{equation}
\end{proposition}

\begin{proof}
The polynomial identity
\begin{equation}
 N^2=F+784C_0\left(Z^2-\frac{X^2}{2}\right)
 \label{eq:theta-conic-identity}
\end{equation}
gives \eqref{eq:theta-cover-identity}, including its scalar, upon
restriction to $D$.  The function-field description in
\eqref{eq:theta-vw} identifies this cover with $C$.

On the conic $C_0=0$, equation \eqref{eq:theta-elimination} becomes
$(X^2-Y^2)^2=0$.  The conics $C_0=0$ and $X^2-Y^2=0$ meet transversely at
the four points $E_{\epsilon,\eta}$.  They have no common point on $Z=0$.
It follows that
\[
 (C_0=0)|_D=2E_0,\qquad
 \operatorname{div}(h)=2E_0-2H.
\]
Taking divisors in \eqref{eq:theta-cover-identity} now gives the asserted
identity of two-torsion classes.  The latter class is nonzero because the
two distinct contact divisors $B_{\ell_+},B_{\ell_-}$ each have
$h^0=1$; if they were linearly equivalent, they would be equal.
\end{proof}

The zeros of $h$ give useful rational charts on the normalized cover.
Near $O_D$, take $X=1+u$, $Z=1$, and solve the quartic equation for $Y$.
Its first terms are
\begin{equation}
 Y=1-u-\frac{48}{49}u^2+O(u^3),\qquad
 h=\frac{u^2}{49}+O(u^3).
 \label{eq:theta-boundary-expansion}
\end{equation}
Indeed substitution of $Y=1-u+au^2+O(u^3)$ gives the coefficient
$-784a-768$ of $u^2$ in $F$.  In the normalized chart $v=u\omega$, the
two points above $O_D$ have $\omega=\pm1/7$.  The same conclusion at the
other three points follows by the sign changes of $X,Y$.  This also
exhibits directly the eight rational boundary points of the cover.

\subsection{Local survival at every finite descent level}

The rational boundary and the real topology have a consequence independent
of the explicit bitangent calculation.  We formulate it for an arbitrary
abelian variety, so it applies both to $\Jac(C)$ and to the quartic quotient.

\begin{theorem}
\label{thm:theta-local-survival}
Let $X/\Q$ be a smooth projective geometrically integral curve, let
$B\subset X$ be a finite closed subset, and suppose $O\in B(\Q)$.
Let $A/\Q$ be an abelian variety and $a:X\to A$ a morphism with $a(O)=0$.
Let $\mathcal V\subset (X\setminus B)(\R)$ contain a point joined to $O$
by a continuous path in $X(\R)$.
For every integer $m\ge1$, form
\[
 X_m=X\times_{a,A,[m]}A.
\]
The geometrically connected component $X_m^0$ containing $(O,0)$ is
defined over $\Q$.  It has a point above $X\setminus B$ over every
$\Q_p$, and a real point above $\mathcal V$.
\end{theorem}

\begin{proof}
In characteristic zero, multiplication by $m$ on $A$ is finite
\emph{\'etale}.  Consequently $X_m\to X$ is finite \emph{\'etale},
and $X_m$ is smooth and projective.  Its geometric connected components
are disjoint.  The unique one containing the rational point $(O,0)$ is
Galois invariant and therefore descends to $\Q$.

Fix a finite prime $p$.  The $p$-adic inverse function theorem, applied
to the projection at $(O,0)$, identifies a neighborhood of this point
with a neighborhood of $O$ in $X(\Q_p)$.  A local parameter on the smooth
curve takes infinitely many values in every sufficiently small
$p$-adic neighborhood.  We may therefore choose a nearby point outside
the finite set $B$ and lift it to $X_m^0(\Q_p)$.  This argument includes
primes dividing $m$: the differential of $[m]$ is multiplication by the
nonzero element $m$ of the field $\Q_p$.

For the real assertion, let $\gamma$ be a path from $O$ to a point
$P\in\mathcal V$.  The path $a\circ\gamma$ starts at zero and lies in
$A(\R)^0$.  This identity component is a compact connected real Lie
group; its multiplication map $[m]$ is a surjective covering map.
Lift $a\circ\gamma$ starting at zero.  Together with $\gamma$ this gives
a path on $X_m(\R)$ beginning at $(O,0)$, and hence lying in $X_m^0$.
Its endpoint is the required point above $P$.
\end{proof}

\begin{corollary}
\label{cor:theta-surviving-covers}
Put $J_C=\Jac(C)$, $J_D=\Jac(D)$, and
\[
 i_C(P)=[P-O_C],\qquad i_D(S)=[S-O_D].
\]
For every $m\ge1$, the distinguished components of both covers
\begin{equation}
 C\times_{i_C,J_C,[m]}J_C,
 \qquad
 C\times_{i_D\circ\pi,J_D,[m]}J_D
 \label{eq:theta-surviving-covers}
\end{equation}
have nonboundary points at every finite place and a real point whose
image on $C$ belongs to $\mathcal U$ and gives $x\ne y$.
\end{corollary}

\begin{proof}
Take $B=C\cap(T=0)$ and choose the positive square roots in
\eqref{eq:theta-C} as $z$ decreases from $+\infty$ to $1/6$.
All three radicands are positive, since they are
$z^2+50$, $(z+7)^2+2$, and $(z-7)^2+2$.
In projective coordinates this path starts at $O_C$.  At $z=1/6$ its
coordinates are
\[
 P=\frac{\sqrt{1801}}6,\qquad
 Q=\frac{\sqrt{1921}}6,\qquad
 R=\frac{\sqrt{1753}}6.
\]
Here $1<P<8$ and $P+7z<8+7/6<50$, so it belongs to $\mathcal U$.
Formula \eqref{eq:theta-return} gives $y-x=7/300\ne0$.
Apply Theorem~\ref{thm:theta-local-survival} with $a=i_C$ and
$a=i_D\circ\pi$, respectively.
\end{proof}

There is also an elementary uniform description of a finite-prime
square sieve on this source.

\begin{proposition}
\label{prop:theta-finite-primes}
For a fixed $z\in\Q_2$, the three equations
\eqref{eq:theta-C} with $T=1$ have simultaneous roots in $\Q_2$ if and
only if $v_2(z)<0$.  The analogous assertion over $\Q_3$ holds with
$v_3(z)<0$.  For every finite set $S$ of primes, there are infinitely
many rational values of $z$ for which all three radicands are squares
in every $\Q_p$, $p\in S$, and whose positive real roots belong to
$\mathcal U$.
\end{proposition}

\begin{proof}
For integral $z$ at two, $z^2+50$ is $2,3$, or $6$ modulo eight, so
cannot be a square.  Write $z=2^{-k}u$ with $k\ge1$ and $u$ a unit.
For $k\ge2$, division by $z^2$ makes all three radicands units congruent
to one modulo eight.  For $k=1$, multiplication by four gives
\[
 u^2+200,\qquad u^2+28u+204,\qquad u^2-28u+204,
\]
which are again one modulo eight.  The complete square criterion in
$\Q_2$ proves sufficiency.

At three, the residue classes $z=0,1,-1$ respectively make the first,
third, and second radicand a nonsquare modulo three.  If $v_3(z)<0$,
division by $z^2$ gives units congruent to one modulo three.  More
generally, at any odd prime $p$, a negative valuation of $z$ gives this
same unit congruence modulo $p$, so Hensel's lemma supplies all roots.

Set $M=6\prod_{p\in S}p$ and take $z=M^{-k}$ for $k\ge1$.
All the required valuations are negative, while $0<z\le1/6$.
At the real place, the positive roots satisfy $1<P<8$ and
$P+7z<8+7/6<50$.  These distinct rational parameters have all the
asserted properties.
\end{proof}

Proposition~\ref{prop:theta-finite-primes} establishes survival of the
simultaneous square conditions at any specified finite set of places.
Corollary~\ref{cor:theta-surviving-covers} establishes the stronger
statement after imposing a fixed multiplication level on either
Jacobian: a single geometrically connected cover retains nonboundary
points at every finite place and a strict real distance configuration.
For the second cover in \eqref{eq:theta-surviving-covers}, this conclusion
includes the literal quadratic lift identified in
Proposition~\ref{prop:theta-cover-class}.  Thus these finite descent
levels have nonempty local-solubility sets with the complete source
and the strict real inequalities imposed.
